%=========================================================
% APA 7th Edition Manuscript Template (LaTeX)
% Built on the apa7 document class + biblatex (APA style).
% Replace every placeholder marked << ... >> with your content.
%=========================================================
\documentclass[man,floatsintex,longtable]{apa7}
% Document-class options (apa7):
%   man          -- manuscript format (title page, double spacing)
%   jou          -- typeset like a printed journal article
%   doc          -- simple document format
%   floatsintext -- place figures/tables in the text (default: at the end)
%   longtable    -- enable multi-page tables
% Example alternative:
%\documentclass[man,floatsintext,longtable]{apa7}

%---------------------------------------------------------
% PACKAGES
%---------------------------------------------------------
\usepackage[utf8]{inputenc}
\usepackage[english]{babel}
\usepackage{amsmath}
\usepackage{amssymb}
\usepackage{graphicx}
\usepackage{subcaption}
% normalem prevents ulem from turning \emph italics into underlines
\usepackage[normalem]{ulem}
\usepackage{xcolor}
% hidelinks removes the colored boxes around hyperlinks
\hypersetup{hidelinks}
\usepackage{nameref}

% Custom commands
\newcommand{\authcomment}[1]{\textcolor{blue}{[#1]}} % editorial note
\newcommand{\red}[1]{\textcolor{red}{#1}}            % flag text in red

%---------------------------------------------------------
% BIBLIOGRAPHY SETUP
%---------------------------------------------------------
\usepackage{csquotes}
\usepackage[style=apa, backend=biber]{biblatex}

%---------------------------------------------------------
% MATH / TABLE / ALGORITHM SETUP
%---------------------------------------------------------
\usepackage{tabularx}
\usepackage{algorithm}
\usepackage{algpseudocode}
\usepackage{bbm}
\usepackage{amsthm}
\usepackage{accents}
\usepackage{multirow}
%\usepackage{setspace}

% Math operator shortcuts
\DeclareMathOperator{\E}{\mathbb{E}}
\DeclareMathOperator{\R}{\mathbb{R}}
\DeclareMathOperator*{\argmaxA}{arg\,max}
\DeclareMathOperator*{\argminA}{arg\,min}
% Centered fixed-width column type: e.g. C{2cm}
\newcolumntype{C}[1]{>{\centering\let\newline\\\arraybackslash\hspace{0pt}}m{#1}}

% Number subfigures as 1, 2, 3 (instead of a, b, c)
\renewcommand\thesubfigure{\arabic{subfigure}}
\setlength{\bibhang}{0.5in}

\addbibresource{ref.bib}

%---------------------------------------------------------
% TITLE PAGE
%---------------------------------------------------------
\title{<<Full Paper Title>>}
\shorttitle{<<Short Title / Running Head>>}
\authorsnames{<<First A. Author>>, <<Second B. Author>>, <<Third C. Author>>}
\authorsaffiliations

\authornote{
{\footnotesize
\textbf{Correspondence:}
Correspondence concerning this article should be addressed to \addORCIDlink{<<Author Name>>}{0000-0000-0000-0000}, <<Mailing address>>. Email: <<email@institution.edu>>.

\textbf{Disclosures statement:} <<State competing interests, or that there are none to declare.>>

\textbf{Funding:} <<State funding source and grant number, or that no funding was received.>>

% \textbf{Prior Dissemination:} <<Preprint location and/or conference presentation, if applicable.>>
}
}

%---------------------------------------------------------
% ABSTRACT
%---------------------------------------------------------
\abstract{<<Abstract (typically 150--250 words): briefly state the problem, approach, key results, and conclusion.>>}

\keywords{<<keyword one, keyword two, keyword three, keyword four>>}

%---------------------------------------------------------
% BEGIN DOCUMENT
%---------------------------------------------------------
\begin{document}
\maketitle

%---------------------------------------------------------
% INTRODUCTION
%---------------------------------------------------------
\section{Introduction}

<<Open with the problem, its significance, and the gap your study addresses. Use \parencite{example-article} for a parenthetical citation and \textcite{example-book} for an in-text citation. End with a roadmap of the paper.>>

\subsection{<<Optional Subsection Heading>>}

<<Subsection content.>>

%---------------------------------------------------------
% METHODS
%---------------------------------------------------------
\section{Methods}

<<Describe the data/sample, materials, design, and analytic procedure in enough detail to be reproducible.>>

% --- Example: a numbered, labeled equation ---
\begin{align}\label{eq:example}
    y = \beta_0 + \beta_1 x + \varepsilon.
\end{align}
Cross-reference it as Equation~\eqref{eq:example}.

% --- Example: a labeled assumptions / items list ---
\begin{itemize}
    \item[(A1)] <<First assumption or item.>>
    \item[(A2)] <<Second assumption or item.>>
\end{itemize}

% --- Example: an algorithm (uncomment to use) ---
% \begin{algorithm}[!h]
% \caption{<<Algorithm name>>}\label{alg:example}
% \begin{algorithmic}[1]
%     \State <<Initialize.>>
%     \For{$i = 1$ to $n$}
%         \State <<Do something.>>
%     \EndFor
%     \State \Return <<output>>
% \end{algorithmic}
% \end{algorithm}

%---------------------------------------------------------
% RESULTS
%---------------------------------------------------------
\section{Results}

<<Report the findings. Reference figures as Figure~\ref{fig:example} and appendix material as Appendix~\ref{app:one} or Table~\ref{app:tab:example}.>>

% --- Example: a single figure ---
% Put image files in the figures/ folder and update the filename below.
\begin{figure}[!h]
    \centering
    \includegraphics[width=1\linewidth]{figures/example-figure.png}
    \caption{<<Figure caption.>>}
    \label{fig:example}
\end{figure}

% --- Example: a figure with side-by-side subfigures (uncomment to use) ---
% \begin{figure}[!h]
%     \centering
%     \begin{subfigure}{0.48\linewidth}
%         \centering
%         \includegraphics[width=\linewidth]{figures/panel-a.png}
%         \caption{<<Panel caption.>>}
%         \label{fig:panel-a}
%     \end{subfigure}
%     \hfill
%     \begin{subfigure}{0.48\linewidth}
%         \centering
%         \includegraphics[width=\linewidth]{figures/panel-b.png}
%         \caption{<<Panel caption.>>}
%         \label{fig:panel-b}
%     \end{subfigure}
%     \caption{<<Overall figure caption.>>}
%     \label{fig:example-multi}
% \end{figure}

%---------------------------------------------------------
% DISCUSSION
%---------------------------------------------------------
\section{Discussion and Conclusions}

<<Interpret the results, connect back to the introduction, acknowledge limitations, and conclude with implications and future directions.>>

%---------------------------------------------------------
% REFERENCES
%---------------------------------------------------------
\printbibliography

%---------------------------------------------------------
% APPENDIX
%---------------------------------------------------------
\appendix
\renewcommand{\thesection}{Appendix \Alph{section}}
% =========================================================
% suppl.tex -- Appendix / Supplementary Material
% Loaded from main.tex via % =========================================================
% suppl.tex -- Appendix / Supplementary Material
% Loaded from main.tex via % =========================================================
% suppl.tex -- Appendix / Supplementary Material
% Loaded from main.tex via \input{suppl} after \appendix.
% Do NOT add \documentclass, \usepackage, or \begin{document} here.
% =========================================================

% Reset counters for the appendix
\setcounter{secnumdepth}{1}
\setcounter{equation}{0}
\setcounter{figure}{0}
\setcounter{table}{0}
%\setcounter{page}{1}   % uncomment to restart page numbering at 1
\setcounter{section}{0}

% Rename counters to S1, S2, ... (supplementary numbering)
\renewcommand{\theequation}{S\arabic{equation}}
\renewcommand{\thefigure}{S\arabic{figure}}
\renewcommand{\thesection}{S\arabic{section}}
\renewcommand{\thetable}{S\arabic{table}}

%---------------------------------------------------------
\section{<<First Appendix Section>>} \label{app:one}

<<Supplementary text. Cross-reference this section from the main text with Appendix~\ref{app:one}.>>

% --- Example: an unnumbered display equation ---
\begin{align*}
    f(x) = \int_{-\infty}^{\infty} g(t)\,dt.
\end{align*}

\newpage

%---------------------------------------------------------
\section{<<Second Appendix Section>>} \label{app:two}

<<Supplementary text.>>

% --- Example: a multi-page table (longtable) ---
\setlength{\tabcolsep}{4pt}          % horizontal padding (default 6pt)
\renewcommand{\arraystretch}{1.0}    % row height
\setlength{\LTpost}{0pt}             % space after the table

\begin{longtable}{@{} l p{0.82\textwidth} @{}}
    \caption{<<Table caption.>>} \label{app:tab:example} \\
    \toprule
    \textbf{<<Column 1>>} & \textbf{<<Column 2>>} \\
    \midrule
    \endfirsthead
    \toprule
    \textbf{<<Column 1>>} & \textbf{<<Column 2>>} \\
    \midrule
    \endhead
    \midrule \multicolumn{2}{r}{\textit{Continued on next page}} \\
    \midrule
    \endfoot
    \bottomrule
    \endlastfoot

    \texttt{<<row key>>} & \small <<row description>> \\
    \texttt{<<row key>>} & \small <<row description>> \\
    \texttt{<<row key>>} & \small <<row description>> \\
\end{longtable}

% Optional notes below the table
\footnotesize
\noindent
\textbf{NOTE}: <<General notes about the table.>>\newline
\textbf{SOURCE}: <<Data source.>>
 after \appendix.
% Do NOT add \documentclass, \usepackage, or \begin{document} here.
% =========================================================

% Reset counters for the appendix
\setcounter{secnumdepth}{1}
\setcounter{equation}{0}
\setcounter{figure}{0}
\setcounter{table}{0}
%\setcounter{page}{1}   % uncomment to restart page numbering at 1
\setcounter{section}{0}

% Rename counters to S1, S2, ... (supplementary numbering)
\renewcommand{\theequation}{S\arabic{equation}}
\renewcommand{\thefigure}{S\arabic{figure}}
\renewcommand{\thesection}{S\arabic{section}}
\renewcommand{\thetable}{S\arabic{table}}

%---------------------------------------------------------
\section{<<First Appendix Section>>} \label{app:one}

<<Supplementary text. Cross-reference this section from the main text with Appendix~\ref{app:one}.>>

% --- Example: an unnumbered display equation ---
\begin{align*}
    f(x) = \int_{-\infty}^{\infty} g(t)\,dt.
\end{align*}

\newpage

%---------------------------------------------------------
\section{<<Second Appendix Section>>} \label{app:two}

<<Supplementary text.>>

% --- Example: a multi-page table (longtable) ---
\setlength{\tabcolsep}{4pt}          % horizontal padding (default 6pt)
\renewcommand{\arraystretch}{1.0}    % row height
\setlength{\LTpost}{0pt}             % space after the table

\begin{longtable}{@{} l p{0.82\textwidth} @{}}
    \caption{<<Table caption.>>} \label{app:tab:example} \\
    \toprule
    \textbf{<<Column 1>>} & \textbf{<<Column 2>>} \\
    \midrule
    \endfirsthead
    \toprule
    \textbf{<<Column 1>>} & \textbf{<<Column 2>>} \\
    \midrule
    \endhead
    \midrule \multicolumn{2}{r}{\textit{Continued on next page}} \\
    \midrule
    \endfoot
    \bottomrule
    \endlastfoot

    \texttt{<<row key>>} & \small <<row description>> \\
    \texttt{<<row key>>} & \small <<row description>> \\
    \texttt{<<row key>>} & \small <<row description>> \\
\end{longtable}

% Optional notes below the table
\footnotesize
\noindent
\textbf{NOTE}: <<General notes about the table.>>\newline
\textbf{SOURCE}: <<Data source.>>
 after \appendix.
% Do NOT add \documentclass, \usepackage, or \begin{document} here.
% =========================================================

% Reset counters for the appendix
\setcounter{secnumdepth}{1}
\setcounter{equation}{0}
\setcounter{figure}{0}
\setcounter{table}{0}
%\setcounter{page}{1}   % uncomment to restart page numbering at 1
\setcounter{section}{0}

% Rename counters to S1, S2, ... (supplementary numbering)
\renewcommand{\theequation}{S\arabic{equation}}
\renewcommand{\thefigure}{S\arabic{figure}}
\renewcommand{\thesection}{S\arabic{section}}
\renewcommand{\thetable}{S\arabic{table}}

%---------------------------------------------------------
\section{<<First Appendix Section>>} \label{app:one}

<<Supplementary text. Cross-reference this section from the main text with Appendix~\ref{app:one}.>>

% --- Example: an unnumbered display equation ---
\begin{align*}
    f(x) = \int_{-\infty}^{\infty} g(t)\,dt.
\end{align*}

\newpage

%---------------------------------------------------------
\section{<<Second Appendix Section>>} \label{app:two}

<<Supplementary text.>>

% --- Example: a multi-page table (longtable) ---
\setlength{\tabcolsep}{4pt}          % horizontal padding (default 6pt)
\renewcommand{\arraystretch}{1.0}    % row height
\setlength{\LTpost}{0pt}             % space after the table

\begin{longtable}{@{} l p{0.82\textwidth} @{}}
    \caption{<<Table caption.>>} \label{app:tab:example} \\
    \toprule
    \textbf{<<Column 1>>} & \textbf{<<Column 2>>} \\
    \midrule
    \endfirsthead
    \toprule
    \textbf{<<Column 1>>} & \textbf{<<Column 2>>} \\
    \midrule
    \endhead
    \midrule \multicolumn{2}{r}{\textit{Continued on next page}} \\
    \midrule
    \endfoot
    \bottomrule
    \endlastfoot

    \texttt{<<row key>>} & \small <<row description>> \\
    \texttt{<<row key>>} & \small <<row description>> \\
    \texttt{<<row key>>} & \small <<row description>> \\
\end{longtable}

% Optional notes below the table
\footnotesize
\noindent
\textbf{NOTE}: <<General notes about the table.>>\newline
\textbf{SOURCE}: <<Data source.>>
   % supplementary material / appendices

\end{document}
